\phantomsection
\addcontentsline{toc}{section}{Lampiran}
\section*{Lampiran}

\begingroup
\singlespacing
\small
\begin{tabular}{@{} l @{\hspace{0.35cm}} l @{}}
  \textbf{a.\enspace Repository} & : \url{https://github.com/MaheswaraKaindra/IF4031_RMB} \\[2pt]
  \multicolumn{2}{@{}l@{}}{\textbf{b.\enspace Form Penggunaan AI Berdasarkan \href{https://drive.google.com/file/d/1OGBapQ118msBHSdf6nLmUWSazS4c5EoO/view}{\underline{Panduan AI ITB}}}} \\
\end{tabular}

\vspace{0.15cm}
\noindent
{\renewcommand{\arraystretch}{1.05}%
\begin{tabularx}{\linewidth}{@{} >{\raggedright\arraybackslash}p{3.8cm} @{\hspace{0.15cm}} p{0.2cm} X @{}}
  Dengan ini, saya/kami* & : & \\
  Nama/NIM               & : & Nathaniel Jonathan Rusli, Maheswara Bayu Kaindra, Buege Mahara Putra / 13523013, 13523015, 13523037 \\
  Fakultas/Sekolah       & : & Sekolah Teknik Elektro dan Informatika - Komputasi \\
  Kode Mata Kuliah       & : & IF4031 \\
  Nama Mata Kuliah       & : & Arsitektur Aplikasi Terdistribusi \\
  Judul Tugas/Proposal   & : & Tugas Besar IF4031 Arsitektur Aplikasi Terdistribusi Milestone 2 \\
\end{tabularx}}

\vspace{0.15cm}
\noindent
Menyatakan bahwa \sout{saya}/\textbf{kami*} \textbf{menggunakan}/\sout{tidak menggunakan}* AI dalam pengerjaan maupun penyusunan luaran tugas pada mata kuliah yang tertulis di atas.

\vspace{0.1cm}
\noindent
Jika Ya, maka alat AI/kecerdasan buatan yang digunakan adalah:

\vspace{0.08cm}
\noindent
{\footnotesize\renewcommand{\arraystretch}{1.15}%
\begin{tabularx}{\linewidth}{| >{\raggedright\arraybackslash}X | >{\centering\arraybackslash}p{2.1cm} | >{\raggedright\arraybackslash}p{3.5cm} |}
  \hline
  \multicolumn{1}{|c|}{\textbf{Lingkup Pekerjaan}} & \textbf{Digunakan?} & \multicolumn{1}{c|}{\textbf{Tingkat Penggunaan}} \\
  \hline
  \textbf{Pemeriksaan Ejaan:} menggunakan \textit{tools} seperti Grammarly, Paperpal, Deepl, Quillbot, Grammarbot, LanguageTool, ProWritingAid, ChatGPT, Google Gemini, atau sejenisnya. & \sout{Ya}/Tidak* & \\
  \hline
  \textbf{Pembuatan Teks:} menggunakan \textit{tools} seperti ChatGPT, Gemini, Paperpal, Copilot, GrammarlyGO, WordAI, WriteSonic, Jasper, Jenni AI, atau sejenisnya. & \sout{Ya}/Tidak* & \\
  \hline
  \textbf{Bantuan Diskusi dalam Penyusunan Konten:} menggunakan \textit{tools} seperti ChatGPT, Gemini, Copilot, Perplexity, Jenni AI, Zoom-Companion, atau sejenisnya. & \sout{Ya}/Tidak* & \\
  \hline
  \textbf{Pembuatan Gambar, Video, dan/atau Grafik:} menggunakan \textit{tools} seperti Craiyon, DALL-E, Midjourney, Stable Diffusion, Microsoft Designer, Gemini, Canva AI, atau sejenisnya. & \sout{Ya}/Tidak* & \\
  \hline
  \textbf{Penulisan Kode Program:} menggunakan AI \textit{coding agent} seperti Claude Code. & Ya/\sout{Tidak}* & \textit{AI task completion-human evaluation} \\
  \hline
\end{tabularx}}

\vspace{0.15cm}
\noindent
\sout{Saya}/\textbf{Kami*} juga menyatakan bahwa setiap penggunaan AI/kecerdasan buatan yang dilakukan pada mata kuliah tersebut telah dinyatakan di dalam surat ini.

\vspace{0.12cm}
\noindent
\hfill Bandung, 22 November 2026

\vspace{0.1cm}
\noindent
\newcommand{\slottd}[1]{%
  \IfFileExists{figures/#1}{%
    \includegraphics[height=1.15cm]{#1}%
  }{%
    \fcolorbox{codeborder}{hintbg}{\parbox[c][0.95cm][c]{3.2cm}{\centering\scriptsize\textcolor{hintgray}{\textit{[Tanda tangan]}}}}%
  }%
}
\begin{tabularx}{\linewidth}{@{} >{\centering\arraybackslash}X >{\centering\arraybackslash}X >{\centering\arraybackslash}X @{}}
  \slottd{ttd-nathan.png} & \slottd{ttd-maheswara.png} & \slottd{ttd-buege.png} \\[3pt]
  Nathaniel Jonathan Rusli & Maheswara Bayu Kaindra & Buege Mahara Putra \\
  NIM 13523013 & NIM 13523015 & NIM 13523037 \\
\end{tabularx}
\endgroup
